% !TeX root = ../../ONMIPA Alin.tex
% Acuan: Zhang (2009), Bab 5, definisi serta Soal 5.1--5.21;
% Zhang (2011), bagian 1.4. Konvensi di sini linear pada argumen kedua,
% sesuai x^*y pada bab matriks; semua rumus konjugasi disesuaikan.
\chapter{Ruang Hasil Kali Dalam}
\label{ch:hasil-kali-dalam}

Dimensi dan kebebasan linear tidak menentukan panjang suatu vektor
atau sudut antarvektor. Kedua gagasan itu memerlukan struktur
tambahan berupa hasil kali dalam. Struktur yang sama juga berlaku
pada ruang polinom dan ruang matriks, sehingga argumen geometris
dapat digunakan untuk menyelesaikan persoalan yang tampaknya aljabar.
Sepanjang bagian ini, medan skalar adalah $\F=\R$ atau $\C$.

\section{Definisi dan Konvensi Konjugasi}

\begin{definisi}[][Hasil kali dalam][hkd:def]
    Hasil kali dalam pada ruang linear $V$ adalah pemetaan
    $\langle\cdot,\cdot\rangle:V\times V\to\F$ yang memenuhi
    \begin{enumerate}[label=(\roman*),nosep]
        \item $\langle x,x\rangle$ real dan nonnegatif, serta
        $\langle x,x\rangle=0$ tepat ketika $x=0$;
        \item $\langle x,y\rangle=\overline{\langle y,x\rangle}$;
        \item $\langle x,ay+bz\rangle
        =a\langle x,y\rangle+b\langle x,z\rangle$.
    \end{enumerate}
    Syarat-syarat tersebut berlaku untuk semua vektor dan skalar
    yang sesuai. Ruang $V$ bersama hasil kali dalamnya disebut
    \emph{ruang hasil kali dalam}.
\end{definisi}

Konvensi handout ini adalah \textbf{linear pada argumen kedua}.
Simetri konjugat kemudian memberi
\[
    \langle ax+by,z\rangle
      =\overline a\langle x,z\rangle+\overline b\langle y,z\rangle.
\]
Jadi pada ruang kompleks, hasil kali dalam bersifat konjugat-linear
pada argumen pertama. Secara umum,
\[
    \left\langle\sum_i a_ix_i,\sum_j b_jy_j\right\rangle
      =\sum_{i,j}\overline{a_i}b_j\langle x_i,y_j\rangle.
\]
Pada ruang real, konjugasi tidak mengubah skalar sehingga bentuknya
bilinear dan simetris. Selain itu, $\langle0,x\rangle
=\langle x,0\rangle=0$.

Hasil kali dalam standar pada $\F^n$ adalah
$\langle x,y\rangle=x^*y=\sum_j\overline{x_j}y_j$.
Untuk vektor real, rumus ini menjadi $x^Ty$.

\begin{exam}[Konjugasi tidak boleh dihilangkan]
    Apakah $B(x,y)=x_1y_1+x_2y_2$ merupakan hasil kali dalam
    pada $\C^2$?
\end{exam}
\begin{solution}
    Tidak. Vektor tak nol $x=(1,i)^T$ memenuhi
    $B(x,x)=1+i^2=0$, sehingga syarat positif definit gagal.
    Bahkan $B((i,0)^T,(i,0)^T)=-1$.
    Pada hasil kali dalam standar, kedua norma kuadrat itu
    masing-masing bernilai dua dan satu.
\end{solution}

Jika $U\subseteq V$ subruang, pembatasan hasil kali dalam pada
$U\times U$ kembali menjadi hasil kali dalam. Namun, suatu
bentuk yang hanya nonnegatif dapat mempunyai vektor tak nol
dengan nilai kuadrat nol; bentuk seperti itu belum menjadi hasil kali dalam.

\section{Hasil Kali Dalam pada Polinom dan Matriks}

\begin{exam}[Hasil kali dalam integral]
    Pada $\mathcal P_{\le n}(\F)$, buktikan bahwa
    \[
        \langle p,q\rangle=\int_{-1}^{1}\overline{p(t)}q(t)\,dt
    \]
    merupakan hasil kali dalam.
\end{exam}
\begin{solution}
    Linearitas dan simetri konjugat mengikuti sifat integral.
    Nilai $\langle p,p\rangle=\int_{-1}^1|p(t)|^2\,dt$ nonnegatif.
    Jika integralnya nol, fungsi kontinu nonnegatif $|p|^2$
    harus nol di seluruh interval: nilai positif di satu titik
    akan memberikan integral positif pada interval kecil di sekitarnya.
    Maka polinom $p$ mempunyai tak hingga banyak akar dan identik nol.
    Arah sebaliknya langsung.
\end{solution}

Pada ruang matriks berukuran $m\times n$, hasil kali dalam Frobenius adalah
\[
    \langle A,B\rangle_F=\tr(A^*B)
      =\sum_{i=1}^m\sum_{j=1}^n\overline{a_{ij}}b_{ij}.
\]
Kesamaan terakhir menunjukkan bahwa ini tepat hasil kali dalam
standar pada daftar seluruh entri. Norma yang dihasilkannya adalah
$\|A\|_F=(\sum_{i,j}|a_{ij}|^2)^{1/2}$.
Matriks satuan $E_{ij}$ saling ortogonal dan masing-masing bernorma satu.

\begin{teorema}[][Representasi oleh matriks positif definit][hkd:matriks-metrik]
    Pada $\F^n$, rumus $\langle x,y\rangle_G=x^*Gy$ merupakan
    hasil kali dalam jika dan hanya jika $G$ Hermitian positif definit.
    Setiap hasil kali dalam pada $\F^n$ mempunyai representasi
    tersebut secara tunggal.
\end{teorema}
\begin{proof}
    Jika $G$ Hermitian positif definit, ketiga aksioma segera
    mengikuti perkalian matriks dan $x^*Gx>0$ untuk $x\ne0$.
    Sebaliknya, letakkan $g_{ij}=\langle e_i,e_j\rangle$.
    Ekspansi kedua argumen dalam basis standar memberi
    $\langle x,y\rangle=x^*Gy$.
    Simetri konjugat memberi $G^*=G$, dan aksioma kepositifan
    memberi sifat positif definit. Setiap entri dipaksa oleh
    hasil kali dua vektor basis, sehingga $G$ unik.
\end{proof}

Untuk basis terurut sembarang $\mathcal B=(b_1,\ldots,b_n)$,
matriks $G_{\mathcal B}=(\langle b_i,b_j\rangle)$ memenuhi
\[
    \langle x,y\rangle=[x]_{\mathcal B}^*G_{\mathcal B}[y]_{\mathcal B}.
\]
Jika $P=P_{\mathcal B\leftarrow\mathcal C}$, maka
$G_{\mathcal C}=P^*G_{\mathcal B}P$.
Jadi perubahan basis pada hasil kali dalam menghasilkan
\emph{kongruensi}, bukan rumus keserupaan $P^{-1}GP$.

\begin{exam}[Parameter dalam hasil kali dalam berbobot]
    Untuk $z\in\C$ dan $a\in\R$, kapan matriks
    \[
        G=\begin{pmatrix}1&z\\\overline z&a\end{pmatrix}
    \]
    mendefinisikan hasil kali dalam pada $\C^2$?
\end{exam}
\begin{solution}
    Untuk $x=(x_1,x_2)^T$, lengkapi kuadrat:
    \[
        x^*Gx=|x_1+zx_2|^2+(a-|z|^2)|x_2|^2.
    \]
    Nilainya positif untuk setiap $x\ne0$ tepat ketika $a>|z|^2$.
    Jika $a=|z|^2$, vektor $(-z,1)^T$ menghasilkan nilai nol.
    Jika $a<|z|^2$, vektor yang sama menghasilkan nilai negatif.
\end{solution}

\section{Norma dan Ketaksamaan Cauchy--Schwarz}

\begin{definisi}[][Norma dan jarak][hkd:norma]
    Norma yang diinduksi oleh hasil kali dalam didefinisikan oleh
    $\|x\|=\sqrt{\langle x,x\rangle}$.
    Jarak antara $x$ dan $y$ adalah $d(x,y)=\|x-y\|$.
    Vektor dengan norma satu disebut vektor satuan.
\end{definisi}

\begin{teorema}[][Ketaksamaan Cauchy--Schwarz][hkd:cauchy-schwarz]
    Untuk setiap $x,y$,
    \[
        |\langle x,y\rangle|^2\le\|x\|^2\|y\|^2.
    \]
    Kesamaan berlaku jika dan hanya jika $x,y$ bergantung linear.
\end{teorema}
\begin{proof}
    Kasus $y=0$ langsung. Untuk $y\ne0$, letakkan
    $c=\langle y,x\rangle/\|y\|^2$ dan $r=x-cy$.
    Maka $\langle y,r\rangle=0$ dan
    \[
        0\le\|r\|^2
          =\|x\|^2-\frac{|\langle y,x\rangle|^2}{\|y\|^2}.
    \]
    Ini memberi ketaksamaan yang diminta.
    Kesamaan ekuivalen dengan $r=0$, yakni $x=cy$.
    Dengan memasukkan kembali kasus vektor nol, kondisi itu
    tepat ketergantungan linear kedua vektor.
\end{proof}

\begin{prop}[Ketaksamaan segitiga dan kesamaannya]
    Norma memenuhi $\|ax\|=|a|\|x\|$ serta
    \[
        \|x+y\|\le\|x\|+\|y\|,\qquad
        \bigl|\|x\|-\|y\|\bigr|\le\|x-y\|.
    \]
    Jika $x,y$ tak nol, kesamaan pada ketaksamaan segitiga
    pertama berlaku tepat ketika $x=cy$ untuk suatu real $c>0$.
    Jika salah satu vektor nol, kesamaan juga berlaku.
\end{prop}
\begin{proof}
    Homogenitas norma mengikuti $\langle ax,ax\rangle
    =|a|^2\langle x,x\rangle$.
    Selanjutnya,
    \begin{align*}
        \|x+y\|^2
        &=\|x\|^2+\|y\|^2+2\operatorname{Re}\langle x,y\rangle\\
        &\le\|x\|^2+\|y\|^2+2|\langle x,y\rangle|\\
        &\le(\|x\|+\|y\|)^2.
    \end{align*}
    Untuk vektor tak nol, kedua ketaksamaan menjadi kesamaan
    jika keduanya bergantung linear dan hasil kali dalamnya
    real positif. Ini memaksa kelipatan real positif.
    Ketaksamaan kedua diperoleh dari
    $\|x\|\le\|x-y\|+\|y\|$, lalu menukar $x,y$.
\end{proof}

Ketergantungan linear saja tidak cukup untuk kesamaan segitiga.
Di $\C$, vektor $1$ dan $i$ bergantung linear, tetapi
$|1+i|=\sqrt2<2=|1|+|i|$.

\begin{exam}[Ketaksamaan berbobot dan kondisi kesamaan]
    Untuk $w_j>0$ dan $a_j\in\C$, buktikan
    \[
        \left|\sum_{j=1}^n a_j\right|^2
        \le\left(\sum_{j=1}^n w_j|a_j|^2\right)
            \left(\sum_{j=1}^n\frac1{w_j}\right).
    \]
\end{exam}
\begin{solution}
    Gunakan Cauchy--Schwarz pada
    $x=(1/\sqrt{w_j})_j$ dan $y=(\sqrt{w_j}a_j)_j$.
    Kesamaan berlaku tepat ketika $y=cx$, yaitu
    $a_j=c/w_j$ untuk suatu konstanta kompleks $c$ yang sama.
    Kondisi kesamaan ini penting ketika ketaksamaan digunakan
    untuk menentukan nilai minimum atau maksimum.
\end{solution}

\section{Identitas Polarisasi dan Ortogonalitas}

\begin{teorema}[][Jajargenjang dan polarisasi][hkd:polarisasi]
    Norma dari hasil kali dalam memenuhi
    \[
        \|x+y\|^2+\|x-y\|^2=2\|x\|^2+2\|y\|^2.
    \]
    Pada ruang real,
    $4\langle x,y\rangle=\|x+y\|^2-\|x-y\|^2$.
    Pada ruang kompleks dengan konvensi handout ini,
    \begin{align*}
        4\langle x,y\rangle
        ={}&\|x+y\|^2-\|x-y\|^2\\
            &-i\|x+iy\|^2+i\|x-iy\|^2.
    \end{align*}
\end{teorema}
\begin{proof}
    Ekspansi $\|x\pm y\|^2$ memberikan suku silang
    $\pm2\operatorname{Re}\langle x,y\rangle$.
    Menjumlahkan memberi identitas jajargenjang, sedangkan
    mengurangkan memberi empat kali bagian real.
    Karena $\langle x,iy\rangle=i\langle x,y\rangle$,
    selisih $\|x+iy\|^2-\|x-iy\|^2$ adalah
    $-4\operatorname{Im}\langle x,y\rangle$.
    Gabungkan bagian real dan imajiner untuk memperoleh rumus kompleks.
\end{proof}

Norma menentukan hasil kali dalamnya secara tunggal melalui
polarisasi. Namun, tidak setiap norma berasal dari hasil kali dalam.
Sebagai contoh, norma $\|(x_1,x_2)\|_1=|x_1|+|x_2|$ memberi
\[
    \|e_1+e_2\|_1^2+\|e_1-e_2\|_1^2=8,
    \qquad 2\|e_1\|_1^2+2\|e_2\|_1^2=4,
\]
sehingga gagal memenuhi identitas jajargenjang.

Dua vektor disebut \emph{ortogonal}, ditulis $x\perp y$, jika
$\langle x,y\rangle=0$. Ortogonalitas mengakibatkan identitas
Pythagoras $\|x+y\|^2=\|x\|^2+\|y\|^2$.
Pada ruang real, arah sebaliknya juga benar.

\begin{exam}[Pythagoras belum cukup di ruang kompleks]
    Di $\C$ dengan hasil kali dalam standar, ambil $x=1$ dan $y=i$.
    Maka $\|x+y\|^2=2=\|x\|^2+\|y\|^2$, tetapi
    $\langle x,y\rangle=i\ne0$.
    Jadi identitas Pythagoras hanya memaksa bagian real hasil
    kali dalam bernilai nol. Pada ruang kompleks, syarat yang tepat adalah
    \[
        x\perp y\iff
        \begin{cases}
            \|x+y\|=\|x-y\|,\\
            \|x+iy\|=\|x-iy\|.
        \end{cases}
    \]
    Kedua persamaan menghilangkan bagian real dan imajiner secara terpisah.
\end{exam}

\section{Matriks Gram dan Ketergantungan Linear}

\begin{definisi}[][Matriks Gram][hkd:gram]
    Untuk vektor $v_1,\ldots,v_m$, matriks Gram adalah
    \[
        G=(g_{ij})_{i,j=1}^m,\qquad g_{ij}=\langle v_i,v_j\rangle.
    \]
\end{definisi}

\begin{teorema}[][Kepositifan dan rank matriks Gram][hkd:rank-gram]
    Matriks Gram Hermitian semidefinit positif dan memenuhi
    \[
        \rank G=\dim\operatorname{span}\{v_1,\ldots,v_m\}.
    \]
    Matriks $G$ positif definit tepat ketika daftar vektor itu
    bebas linear; ekuivalen dengan $\det G>0$.
\end{teorema}
\begin{proof}
    Untuk $c=(c_1,\ldots,c_m)^T$, berlaku
    \[
        c^*Gc=\left\|\sum_{j=1}^m c_jv_j\right\|^2\ge0.
    \]
    Simetri konjugat memberi $G^*=G$.
    Definisikan $S:\F^m\to V$ melalui $Sc=\sum c_jv_j$.
    Jika $Gc=0$, persamaan di atas memberi $Sc=0$.
    Sebaliknya, $Sc=0$ memberi $(Gc)_i=\langle v_i,Sc\rangle=0$.
    Maka $\ker G=\ker S$; teorema rank--nulitas memberi
    $\rank G=\rank S=\dim\operatorname{im}S$.
    Kepositifan definit ekuivalen dengan $Sc\ne0$ untuk setiap
    $c\ne0$, yakni kebebasan linear. Karena semua nilai eigen
    $G$ nonnegatif, keadaan ini juga ekuivalen dengan $\det G>0$.
\end{proof}

Dalam koordinat ortonormal, jika kolom matriks $A$ adalah para
$v_j$, maka $G=A^*A$. Sebaliknya, setiap matriks Hermitian
semidefinit positif merupakan matriks Gram: ambil akar kuadrat
Hermitian $B=G^{1/2}$ sehingga $G=B^*B$, lalu gunakan kolom-kolom $B$.
Hasil mengenai akar kuadrat telah dibahas pada bab matriks khusus.

\begin{exam}[Tiga hasil kali dalam tidak dapat dipilih sembarang]
    Misalkan $u,v,w$ vektor satuan dalam ruang real dan
    $a=\langle u,v\rangle$, $b=\langle u,w\rangle$,
    $c=\langle v,w\rangle$. Buktikan
    \[
        (c-ab)^2\le(1-a^2)(1-b^2).
    \]
\end{exam}
\begin{solution}
    Matriks Gram ketiganya semidefinit positif, sehingga
    \[
        0\le\det\begin{pmatrix}1&a&b\\a&1&c\\b&c&1\end{pmatrix}
          =1+2abc-a^2-b^2-c^2.
    \]
    Menata ulang memberi ketaksamaan yang diminta.
    Kesamaan berarti determinan Gram nol, sehingga ketiga
    vektor bergantung linear. Jika $|a|,|b|<1$, batas untuk $c$ adalah
    $ab\pm\sqrt{(1-a^2)(1-b^2)}$.
\end{solution}

\section{Latihan Pilihan}

\begin{latihan}\label{hkd:lat-tarik-balik}
    Misalkan $A:V\to W$ linear dan $W$ ruang hasil kali dalam.
    Buktikan bahwa $[x,y]=\langle Ax,Ay\rangle$ merupakan hasil
    kali dalam pada $V$ tepat ketika $A$ injektif.
    Nyatakan syaratnya jika $A$ adalah matriks $m\times n$.
\end{latihan}

\begin{latihan}\label{hkd:lat-evaluasi}
    Pada $\mathcal P_{\le n}(\F)$, ambil titik real berbeda
    $a_1,\ldots,a_m$ dan bobot $w_j>0$. Tentukan kapan
    \[
        [p,q]=\sum_{j=1}^m w_j\overline{p(a_j)}q(a_j)
    \]
    merupakan hasil kali dalam.
\end{latihan}

\begin{latihan}\label{hkd:lat-turunan}
    Pada $\mathcal P_{\le n}(\F)$, buktikan bahwa
    \[
        [p,q]=\overline{p(0)}q(0)+\int_0^1\overline{p'(t)}q'(t)\,dt
    \]
    merupakan hasil kali dalam. Jelaskan mengapa menghapus
    suku pertama membuat sifat positif definit gagal.
\end{latihan}

\begin{latihan}\label{hkd:lat-kesamaan}
    Di $\C^2$, ambil $x=(1,i)^T$ dan $y=(a,b)^T$.
    Tentukan semua $a,b$ yang memberikan kesamaan pada
    Cauchy--Schwarz, lalu tentukan yang juga memenuhi
    $\|x+y\|=\|x\|+\|y\|$.
\end{latihan}

\begin{latihan}\label{hkd:lat-trace}
    Untuk $A,B\in\Mat_n(\C)$, buktikan
    \[
        |\tr A|^2\le n\tr(A^*A),\qquad
        |\tr(AB)|\le\|A\|_F\|B\|_F.
    \]
    Tentukan kondisi kesamaan pada ketaksamaan pertama.
\end{latihan}

\begin{latihan}\label{hkd:lat-invers-positif}
    Jika $G$ Hermitian positif definit dan $x\ne0$, buktikan
    \[
        (x^*Gx)(x^*G^{-1}x)\ge(x^*x)^2.
    \]
    Buktikan bahwa kesamaan berlaku tepat ketika $x$ vektor eigen $G$.
\end{latihan}

\begin{latihan}\label{hkd:lat-korelasi}
    Untuk $m\ge2$ dan $a\in\R$, tentukan kapan terdapat
    vektor satuan real $v_1,\ldots,v_m$ dengan
    $\langle v_i,v_j\rangle=a$ untuk setiap $i\ne j$.
    Tentukan dimensi terkecil ruang yang dapat memuat vektor-vektor itu.
\end{latihan}

\begin{latihan}\label{hkd:lat-negatif}
    Dalam ruang hasil kali dalam real berdimensi $n\ge1$,
    misalkan $v_1,\ldots,v_m$ memenuhi
    $\langle v_i,v_j\rangle<0$ untuk setiap $i\ne j$.
    Buktikan $m\le n+1$ dan tunjukkan bahwa batas tersebut dapat dicapai.
\end{latihan}

\ifpembahasan
\section{Pembahasan Latihan}

\subsection*{Latihan \ref{hkd:lat-tarik-balik}: kepositifan ditentukan oleh kernel}
Linearitas pada argumen kedua dan simetri konjugat selalu berlaku.
Selanjutnya $[x,x]=\|Ax\|^2=0$ tepat ketika $Ax=0$.
Jadi sifat positif definit ekuivalen dengan $\ker A=\{0\}$.
Untuk matriks $m\times n$, ini berarti $\rank A=n$;
khususnya, diperlukan $m\ge n$.

\subsection*{Latihan \ref{hkd:lat-evaluasi}: banyak titik dibandingkan derajat}
Karena semua bobot positif, $[p,p]=0$ tepat ketika
$p(a_1)=\cdots=p(a_m)=0$.
Jika $m\ge n+1$, polinom berderajat paling besar $n$ dengan
akar sebanyak itu harus nol, sehingga bentuknya hasil kali dalam.
Jika $m\le n$, polinom tak nol $\prod_{j=1}^m(t-a_j)$
berada di ruang yang diberikan dan bernilai nol pada semua titik.
Jadi syaratnya tepat $m\ge n+1$.

\subsection*{Latihan \ref{hkd:lat-turunan}: memulihkan konstanta yang hilang}
Jika $[p,p]=0$, kedua suku nonnegatifnya nol.
Integral $|p'|^2$ yang nol memaksa $p'=0$, sehingga $p$ konstan.
Syarat $p(0)=0$ kemudian memaksa konstanta itu nol.
Aksioma lainnya mengikuti sifat konjugasi dan integral.
Tanpa suku pertama, setiap polinom konstan tak nol mempunyai
nilai kuadrat nol, sehingga bentuknya bukan hasil kali dalam.

\subsection*{Latihan \ref{hkd:lat-kesamaan}: bergantung linear dan searah}
Karena $x\ne0$, kesamaan Cauchy--Schwarz berarti $y=cx$,
yakni $a=c$ dan $b=ic$ untuk suatu $c\in\C$.
Kesamaan segitiga mengharuskan $c$ real nonnegatif, termasuk
$c=0$ yang menghasilkan $y=0$.
Kelipatan kompleks dengan fase berbeda umumnya tidak memenuhi
kesamaan segitiga.

\subsection*{Latihan \ref{hkd:lat-trace}: pandang matriks sebagai vektor}
Gunakan $\langle I,A\rangle_F=\tr A$ dan $\|I\|_F^2=n$.
Cauchy--Schwarz memberi ketaksamaan pertama, dengan kesamaan
tepat ketika $A=cI$ untuk suatu $c\in\C$.
Untuk yang kedua, gunakan
$\langle A^*,B\rangle_F=\tr(AB)$ serta $\|A^*\|_F=\|A\|_F$.

\subsection*{Latihan \ref{hkd:lat-invers-positif}: gunakan metrik yang sesuai}
Gunakan Cauchy--Schwarz pada $x,G^{-1}x$ untuk hasil kali dalam
$\langle u,v\rangle_G=u^*Gv$. Ketiga besaran yang muncul adalah
\[
    \langle x,G^{-1}x\rangle_G=x^*x,\quad
    \|x\|_G^2=x^*Gx,\quad
    \|G^{-1}x\|_G^2=x^*G^{-1}x.
\]
Kesamaan berarti $G^{-1}x=cx$, dengan $c>0$ karena $G^{-1}$
positif definit. Ini ekuivalen dengan $x$ vektor eigen $G$.

\subsection*{Latihan \ref{hkd:lat-korelasi}: spektrum matriks Gram}
Matriks Gram yang diminta adalah
$G=(1-a)I_m+a\boldsymbol{1}\boldsymbol{1}^T$.
Pada arah $\boldsymbol{1}$, nilai eigennya $1+(m-1)a$;
pada komplemen ortogonalnya, nilai eigennya $1-a$.
Jadi $G$ semidefinit positif tepat ketika
\[
    -\frac1{m-1}\le a\le1.
\]
Kondisi ini cukup karena setiap matriks real simetris semidefinit
positif dapat diwujudkan sebagai Gram vektor real.
Dimensi terkecil sama dengan rank Gram: bernilai $m$ di bagian
dalam interval, $m-1$ pada ujung kiri, dan $1$ pada ujung kanan.
Kecukupan dimensi rank diperoleh dengan mengambil hanya baris
tak nol pada faktorisasi spektral $G=B^TB$.

\subsection*{Latihan \ref{hkd:lat-negatif}: pisahkan koefisien positif dan negatif}
Andaikan $m>n+1$. Vektor $(v_j,1)$ pada $V\times\R$
bergantung linear, sehingga ada real $c_j$ tidak semuanya nol dengan
\[
    \sum_jc_jv_j=0,\qquad \sum_jc_j=0.
\]
Ada koefisien positif dan negatif. Letakkan
$u=\sum_{c_j>0}c_jv_j=\sum_{c_k<0}(-c_k)v_k$.
Menggunakan dua ekspresi tersebut pada kedua argumen memberi
\[
    \|u\|^2=\sum_{c_j>0,\ c_k<0}c_j(-c_k)\langle v_j,v_k\rangle<0,
\]
kontradiksi. Jadi $m\le n+1$.
Untuk mencapai batas, di hiperbidang $\sum x_j=0$ pada $\R^{n+1}$
yang berdimensi $n$, ambil
$v_j=e_j-\boldsymbol{1}/(n+1)$.
Untuk $j\ne k$, hasil kali dalamnya $-1/(n+1)<0$.
\fi
